\documentclass[11pt]{article}
\usepackage[margin=1.2in]{geometry}
\usepackage{amsmath,amssymb,amsthm}
\usepackage[hidelinks]{hyperref}
\hypersetup{
  pdftitle={Nonexistence of consecutive powerful triplets around cubes with mixed prime factorizations},
  pdfauthor={Berkay Y\"uksel Sayim (ORCID 0009-0004-4993-7352)},
  pdfsubject={Number theory: the Erdos-Mollin-Walsh conjecture},
  pdfkeywords={powerful numbers, consecutive powerful numbers, Erdos-Mollin-Walsh conjecture, Mordell equation, genus 2 curve, elliptic curve, cyclotomic polynomial}
}

\newtheorem{theorem}{Theorem}
\newtheorem{lemma}[theorem]{Lemma}
\newtheorem{proposition}[theorem]{Proposition}
\newtheorem{corollary}[theorem]{Corollary}
\theoremstyle{remark}
\newtheorem{remark}[theorem]{Remark}

\newcommand{\Z}{\mathbb{Z}}
\newcommand{\Q}{\mathbb{Q}}
\newcommand{\F}{\mathbb{F}}
\newcommand{\vp}{v}

\title{Nonexistence of consecutive powerful triplets around cubes\\ with mixed prime factorizations}
\author{Berkay Y\"uksel Sayim\\ \small Independent Researcher, Germany\\ \small \texttt{berksa@tutamail.com} ~$\cdot$~ ORCID 0009-0004-4993-7352}
\date{June 2026}

\begin{document}
\maketitle

\begin{abstract}
A positive integer is \emph{powerful} if every prime in its factorization occurs with exponent at least two. The Erd\H{o}s--Mollin--Walsh conjecture asserts that no three consecutive integers are all powerful. Chan (2025) excluded triples of the form $x^3-1=p^3y^2,\; x^3,\; x^3+1=q^3z^2$, and She (2025) excluded $x^3-1=p^2a^3,\; x^3,\; x^3+1=q^2b^3$, with $p,q$ prime. We close the remaining \emph{mixed} configurations: there are no consecutive powerful triples of the form $x^3-1=p^2a^3,\ x^3+1=q^3z^2$, nor of the form $x^3-1=p^3y^2,\ x^3+1=q^2b^3$. The case analysis reduces both families to the single Diophantine equation $t^6+t^3+1=3w^2$, which we solve completely: exploiting the palindromic symmetry of the genus-$2$ curve $3w^2=t^6+t^3+1$, we descend to an elliptic quotient of Mordell--Weil rank $0$ (curve \texttt{1296.k1} in the LMFDB, Cremona label \texttt{1296f2}), so the only rational solutions are $(t,w)=(1,\pm 1)$. We further show that no finite congruence sieve can decide this equation, explaining the necessity of the elliptic argument. As a complementary result we refine Beckon's mod-$36$ constraint on the smallest member of a hypothetical triple to arbitrary prime-square moduli and quantify the density decay of admissible residue classes.
\end{abstract}

\section{Introduction}

A positive integer $m$ is \emph{powerful} if $p^2 \mid m$ for every prime $p \mid m$; equivalently $m=a^2b^3$ with $a,b\in\Z_{>0}$ \cite{Golomb}. Pairs of consecutive powerful numbers, such as $(8,9)$ and $(288,289)$, occur infinitely often (e.g.\ via the Pell equation $x^2-8y^2=1$). In contrast, Erd\H{o}s \cite{Erdos76} conjectured, and Mollin and Walsh \cite{MollinWalsh86a,MollinWalsh86b} investigated systematically, that \emph{no three consecutive integers are all powerful}. The conjecture remains open; under the $abc$-conjecture at most finitely many triples exist.

If $(n,n+1,n+2)$ is such a triple, the middle member may or may not be a perfect cube. The case in which the middle member is a cube $x^3$ has attracted recent attention. Chan \cite{Chan2025} proved that there is no triple of the shape
\[
x^3-1=p^3y^2,\qquad x^3,\qquad x^3+1=q^3z^2 \qquad (p,q \text{ prime}),
\]
and She \cite{She2025} proved the analogous statement for
\[
x^3-1=p^2a^3,\qquad x^3,\qquad x^3+1=q^2b^3 \qquad (p,q \text{ prime}).
\]
Both works treat \emph{symmetric} configurations: the two outer members carry the same shape ($p^3\cdot\square$ on both sides in \cite{Chan2025}; $p^2\cdot\text{cube}$ on both sides in \cite{She2025}). Further recent activity in this circle includes van Doorn's study of three-term arithmetic progressions of consecutive powerful numbers \cite{vanDoorn}. The natural \emph{mixed} configurations have, to our knowledge, not been treated. We close them.

\begin{theorem}\label{thm:main}
There are no integers $x\ge 2$, primes $p,q$, and integers $a,b,y,z$ such that either
\begin{align}
\mathrm{(M1)}\qquad & x^3-1=p^2a^3 \quad\text{and}\quad x^3+1=q^3z^2, \label{eq:M1}\\
\mathrm{(M2)}\qquad & x^3-1=p^3y^2 \quad\text{and}\quad x^3+1=q^2b^3. \label{eq:M2}
\end{align}
In particular, no triple of consecutive powerful numbers $(x^3-1,\,x^3,\,x^3+1)$ has mixed outer factorizations of these shapes.
\end{theorem}

The proof is a case analysis in the spirit of \cite{Chan2025,She2025} (Sections~\ref{sec:M1}--\ref{sec:M2}): factoring $x^3\mp 1=(x\mp 1)(x^2\pm x+1)$ and tracking $3$-adic valuations reduces all but one branch of each family to classical Mordell equations $Y^2=X^3+k$ with $k\in\{2,-2,54,-54,-162\}$, whose integral points are completely known \cite{GPZ,BennettGhadermarzi,LMFDB}. Both remaining branches lead to one and the same equation,
\begin{equation}\label{eq:wall}
t^6+t^3+1=3w^2,
\end{equation}
with $t\ge 7$ in family (M1) and $t\le -5$ in family (M2). Equation~\eqref{eq:wall} concerns the ninth cyclotomic value $\Phi_9(t)$ and does not yield to the degree-$4$ or congruence techniques of \cite{Chan2025,She2025}; indeed we prove (Remark~\ref{rem:covering}) that \emph{no} finite congruence sieve can decide it. Instead we solve it completely:

\begin{theorem}\label{thm:wall}
The only solutions of $t^6+t^3+1=3w^2$ in rational numbers are $(t,w)=(1,\pm 1)$.
\end{theorem}

The proof (Section~\ref{sec:wall}) exploits the palindromic symmetry $t\mapsto 1/t$ of the genus-$2$ curve $3w^2=t^6+t^3+1$: the curve admits two independent elliptic quotients over $\Q$, and one of them has Mordell--Weil rank $0$ with trivial torsion --- it is the curve \texttt{1296.k1} in the LMFDB \cite{LMFDB}, Cremona label \texttt{1296f2} --- which pins down all rational points. Since $t=1$ lies in neither branch range, both remaining branches are empty and Theorem~\ref{thm:main} follows.

As a complementary result in the same circle of ideas we record a modular refinement of a little-known constraint of Beckon \cite{Beckon} on arbitrary consecutive powerful triples (not assuming a cube in the middle).

\begin{proposition}\label{prop:beckon}
Let $(n,n+1,n+2)$ be powerful. Then, beyond Beckon's constraint $n \bmod 36\in\{7,27,35\}$:
\begin{enumerate}
\item[(i)] $n \bmod 900$ lies in an explicit set of $39$ residue classes, and $n \bmod 44100$ in an explicit set of $1209$ classes;
\item[(ii)] for every prime $\ell \ge 5$, the proportion of admissible residues modulo $\ell^2$ is $(\ell^2-3\ell+3)/\ell^2$, while it is $1/4$ for $\ell=2$ and $1/3$ for $\ell=3$;
\item[(iii)] consequently the density of admissible classes modulo $\prod_{\ell\le L}\ell^2$ tends to $0$ as $L\to\infty$, but only at the rate $\asymp (\log L)^{-3}$; in particular no finite sieve of this type can resolve the Erd\H{o}s--Mollin--Walsh conjecture.
\end{enumerate}
\end{proposition}

All computer calculations below are short and reproducible; the elliptic-curve computations were performed in SageMath and independently cross-checked (two independent runs, plus the LMFDB data \cite{LMFDB}); the symbolic identities were verified exactly in SymPy. The verification scripts are provided as ancillary files. 

\section{Preliminaries}\label{sec:prelim}

Throughout, $x\ge 2$ is an integer, $p,q$ are primes, and $\vp_\ell$ denotes the $\ell$-adic valuation. We collect standard facts; proofs are short and included for completeness where not cited.

\begin{lemma}[parity]\label{lem:parity}
In any triple of consecutive powerful numbers the even member is divisible by $4$. If the middle member is $x^3$, then $x$ is even; in particular $x^3\pm1$ are odd and coprime, so $p\neq q$ in \eqref{eq:M1}--\eqref{eq:M2}.
\end{lemma}
\begin{proof}
Exactly one or two of three consecutive integers are even; a powerful even number is $\equiv 0 \pmod 4$. If $n$ and $n+2$ were both even, one of them would be $\equiv 2\pmod 4$, hence not powerful. So the even member is unique, namely the middle one if the outer two are odd. For the middle member $x^3$: if $x$ were odd, then $x^3\pm1$ would be even, contradiction. Finally $\gcd(x^3-1,x^3+1)\mid 2$ and both are odd.
\end{proof}

\begin{lemma}[gcd and LTE]\label{lem:gcd}
$x^3\mp1=(x\mp1)(x^2\pm x+1)$ with $g_\mp:=\gcd(x\mp 1,\,x^2\pm x+1)=\gcd(x\mp1,3)\in\{1,3\}$. Moreover, if $3\mid x\mp 1$ then $\vp_3(x^2\pm x+1)=1$ exactly.
\end{lemma}
\begin{proof}
$x^2\pm x+1=(x\mp1)(x\pm2)+3$, giving the gcd claim (cf.\ \cite[Lemma~4]{She2025}). For the valuation: if $x=1+3k$ then $x^2+x+1=3+9k+9k^2\equiv3\pmod9$; if $x=-1+3k$ then $x^2-x+1=3-9k+9k^2\equiv3\pmod9$. In both cases $\vp_3=1$ exactly.
\end{proof}

\begin{lemma}[square sandwich]\label{lem:sandwich}
For $x\ge2$: $x^2<x^2+x+1<(x+1)^2$ and $(x-1)^2<x^2-x+1<x^2$. Hence neither $x^2+x+1$ nor $x^2-x+1$ is a perfect square.
\end{lemma}

\begin{lemma}[Chan-side structure]\label{lem:chanside}
Suppose $x^3+1=q^3z^2$ and $g_+=1$. Then $x+1$ is a perfect square and $x^2-x+1=q^{3+2k}\cdot\square$ for some $k\ge0$. Analogously, if $x^3-1=p^3y^2$ and $g_-=1$, then $x-1$ is a perfect square and $x^2+x+1=p^{3+2k}\cdot\square$.
\end{lemma}
\begin{proof}
With $g_+=1$ the factors $x+1$ and $x^2-x+1$ are coprime, so each is, up to the prime $q$, a perfect square, and $q^3$ (odd exponent $\ge3$) divides exactly one of them. If $q^3$ divided $x+1$ with $x^2-x+1$ a square, Lemma~\ref{lem:sandwich} is contradicted. The mirrored statement is identical.
\end{proof}

\begin{lemma}[She's splitting; {\cite[Lemma~1]{She2025}}]\label{lem:she1}
If $RS=p^2C^3$ with $\gcd(R,S)=1$, then one of $R,S$ is a perfect cube and the other is $p^2$ times a perfect cube.
\end{lemma}

\begin{lemma}[{\cite[Lemma~2]{She2025}}]\label{lem:she2}
$u^2+u+1=3v^3$ has only the integer solutions $u\in\{-2,1\}$, and $u^2-u+1=3v^3$ only $u\in\{-1,2\}$.
\end{lemma}

\begin{lemma}[Tzanakis; {\cite{Tzanakis}, cf.\ \cite[Lemma~3]{She2025}}]\label{lem:she3}
$u^2+u+1=v^3$ has only $u\in\{-19,-1,0,18\}$, and $u^2-u+1=v^3$ only $u\in\{-18,0,1,19\}$.
\end{lemma}

\begin{lemma}[Mordell curves]\label{lem:mordell}
The integral points of $Y^2=X^3+k$ are: for $k=2$: $(-1,\pm1)$; for $k=-2$: $(3,\pm5)$; for $k=54$: $(3,\pm9)$; for $k=-54$: $(7,\pm17)$; for $k=-162$: none.
\end{lemma}
\begin{proof}
All five values satisfy $|k|\le 10^4$ and are covered by the complete solution of Mordell's equation in this range \cite{GPZ}, extended to $|k|\le 10^7$ in \cite{BennettGhadermarzi}; see also the corresponding LMFDB data \cite{LMFDB}. (The case $k=-2$ is classical.) We re-verified all five lists with \texttt{integral\_points} in SageMath in two independent runs. 
\end{proof}

We will also use without further comment: squares mod $9$ are $\{0,1,4,7\}$; cubes mod $9$ are $\{0,1,8\}$; cubes of integers $\equiv1\pmod3$ are $\equiv\{1,10,19\}\pmod{27}$.

\paragraph{Exponent bookkeeping.} For the Chan shape $x^3\mp1=P^3Y^2$ one has $\vp_P\equiv1\pmod2$, $\vp_P\ge3$, and $\vp_r\equiv0\pmod 2$ for primes $r\neq P$; for the She shape $x^3\mp1=P^2A^3$ one has $\vp_P\equiv2\pmod3$ and $\vp_r\equiv0\pmod3$ for $r\neq P$. Since $g_\mp\in\{1,3\}$, every prime $r\neq3$ divides exactly one of the two factors $x\mp1$, $x^2\pm x+1$.

\section{Family (M1): $x^3-1=p^2a^3$, $x^3+1=q^3z^2$}\label{sec:M1}

By Lemma~\ref{lem:parity}, $x$ is even and $p\neq q$. We split by $x\bmod 3$.

\subsection*{Case 1: $x\equiv0\pmod3$ (so $g_-=g_+=1$)}
By Lemma~\ref{lem:chanside}, $x+1=v^2$. By Lemma~\ref{lem:she1} applied to $(x-1)(x^2+x+1)=p^2a^3$:
\begin{itemize}
\item[(ii)] $x-1=p^2u^3$, $x^2+x+1=t^3$: Lemma~\ref{lem:she3} forces $x\in\{-19,-1,0,18\}$; only $x=18$ is even and $\ge2$, but $x-1=17$ is prime, not of the form $p^2u^3$. \emph{Closed.}
\item[(i)] $x-1=u^3$, $x^2+x+1=p^2t^3$: combining with $x+1=v^2$ gives $v^2-u^3=2$, a point on $Y^2=X^3+2$; by Lemma~\ref{lem:mordell} $u=-1$, $x=0<2$. \emph{Closed.}
\end{itemize}

\subsection*{Case 2: $x\equiv1\pmod3$ (so $g_-=3$, $g_+=1$)}
Here $x^2-x+1\equiv1\pmod3$, so $q\neq3$, and $x+1=v^2$ by Lemma~\ref{lem:chanside}; on the She side $\vp_3(x^2+x+1)=1$ exactly (Lemma~\ref{lem:gcd}).
\begin{itemize}
\item $p=3$: here $\vp_3(x^3-1)=2+3\vp_3(a)$, while $\vp_3(x^3-1)=\vp_3(x-1)+1$ by Lemma~\ref{lem:gcd}; hence $\vp_3(x-1)=1+3\vp_3(a)$ and $\vp_3(x^2+x+1)=1$. All prime exponents in $x^2+x+1$ other than that of $3$ are $\equiv0\pmod3$, so $x^2+x+1=3t^3$ with $3\nmid t$; Lemma~\ref{lem:she2} gives $x\in\{-2,1\}$, both $<2$. \emph{Closed.}
\item $p\neq3$, $p\mid x-1$: then again $x^2+x+1=3t^3$ ($\vp_3=1$, $p\nmid$, all other exponents $\equiv0\bmod3$), so $x\in\{-2,1\}$. \emph{Closed.}
\item $p\neq3$, $p\mid x^2+x+1$: then $\vp_3(x-1)=3\vp_3(a)-1\equiv2\pmod3$ and all other prime exponents in $x-1$ are $\equiv0\pmod3$, hence $x-1=9d^3$. But $x+1=v^2$ gives $v^2=9d^3+2\equiv2\pmod 9$, impossible since squares mod $9$ are $\{0,1,4,7\}$. \emph{Closed.}
\end{itemize}

\subsection*{Case 3: $x\equiv2\pmod3$ (so $g_-=1$, $g_+=3$)}
She side (Lemma~\ref{lem:she1}):
\begin{itemize}
\item[(ii)] $x-1=p^2u^3$, $x^2+x+1=t^3$: Lemma~\ref{lem:she3} gives $x\in\{-19,-1,0,18\}$; none is $\equiv2\pmod 3$, even, and $\ge2$. \emph{Closed.}
\item[(i)] $x-1=u^3$ with $u$ odd, $u\equiv1\pmod3$ (so $u\equiv1\bmod 6$), and $x^2+x+1=p^2t^3$. On the Chan side $g_+=3$ and $\vp_3(x^2-x+1)=1$ exactly. Sub-split on $q$:
  \begin{itemize}
  \item $q=3$: $x^3+1=27z^2$ gives $\vp_3(x+1)=2+2\vp_3(z)$ even and all other exponents in $x+1$ even, so $x+1=m^2$; then $m^2-u^3=2$, again $Y^2=X^3+2$, so $u=-1$, $x=0<2$. \emph{Closed.}
  \item $q\neq3$, $q^3\mid x^2-x+1$: then $\vp_3(x+1)$ is odd, say $x+1=3^{2j+1}m^2$ with $\gcd(m,3)=1$ up to squares. If $j\ge1$ then $27\mid x+1=u^3+2$, i.e.\ $u^3\equiv-2\equiv25\pmod{27}$, impossible for $u\equiv1\pmod3$. So $j=0$, $x+1=3m^2$, i.e.\ $u^3=3m^2-2$; multiplying by $27$: $(9m)^2=(3u)^3+54$. By Lemma~\ref{lem:mordell} the only points have $X=3u=3$, so $u=1$, $x=2$; but $x^3-1=7$ is not powerful. \emph{Closed.}
  \item $q\neq3$, $q^3\mid x+1$: then $x^2-x+1=3w^2$ ($\vp_3=1$ exactly, $q\nmid$, all other exponents even). With $x=u^3+1$:
  \[
  (u^3+1)^2-(u^3+1)+1=u^6+u^3+1=3w^2,
  \]
  i.e.\ equation~\eqref{eq:wall} with $t=u$. Here $u\ge1$ since $x\ge2$, and $u=1$ gives $x=2$, where $x^3-1=7$ is not of the form $p^2a^3$; together with $u\equiv1\pmod 6$ this forces $t=u\ge7$. \emph{This is the only surviving branch of (M1); it is empty by Theorem~\ref{thm:wall}.}
  \end{itemize}
\end{itemize}

\section{Family (M2): $x^3-1=p^3y^2$, $x^3+1=q^2b^3$}\label{sec:M2}

Again $x$ is even, $p\neq q$; split by $x\bmod3$.

\subsection*{Case 1: $x\equiv0\pmod3$ ($g_-=g_+=1$)}
Chan side (Lemma~\ref{lem:chanside}): $x-1=u^2$ with $u$ odd. She side (Lemma~\ref{lem:she1}) applied to $(x+1)(x^2-x+1)=q^2b^3$:
\begin{itemize}
\item[(b)] $x+1=q^2s^3$, $x^2-x+1=t^3$: Lemma~\ref{lem:she3} gives $x\in\{-18,0,1,19\}$; no admissible value. \emph{Closed.}
\item[(a)] $x+1=s^3$, $x^2-x+1=q^2t^3$: then $s^3-u^2=2$, i.e.\ a point on $Y^2=X^3-2$; by Lemma~\ref{lem:mordell} $(s,u)=(3,\pm5)$, $x=26\equiv2\pmod 3$, contradicting the case. \emph{Closed.}
\end{itemize}

\subsection*{Case 2: $x\equiv1\pmod3$ ($g_-=3$, $g_+=1$)}
She side has $g_+=1$:
\begin{itemize}
\item[(b)] $x+1=q^2s^3$, $x^2-x+1=t^3$: as above, no admissible $x$. \emph{Closed.}
\item[(a)] $x+1=s^3$ with $s$ odd, $s\equiv2\pmod3$, hence $s\equiv5\pmod 6$ and $s\ge5$; and $x^2-x+1=q^2t^3$. On the Chan side $\vp_3(x^2+x+1)=1$ exactly. Sub-split on $p$:
  \begin{itemize}
  \item $p=3$: then $\vp_3(x-1)$ is even and all other exponents in $x-1$ are even, so $x-1=u^2$; then $u^2=s^3-2$, so $(s,u)=(3,\pm5)$ by Lemma~\ref{lem:mordell}, giving $x=26\equiv2\pmod3\neq1$. \emph{Closed.}
  \item $p\neq3$, $p\mid x-1$: then $x^2+x+1=3w^2$ ($\vp_3=1$, $p\nmid$, rest even). With $x=s^3-1$:
  \[
  (s^3-1)^2+(s^3-1)+1=s^6-s^3+1=3w^2,
  \]
  i.e.\ equation~\eqref{eq:wall} with $t=-s\le-5$. \emph{This is the only surviving branch of (M2); it is empty by Theorem~\ref{thm:wall}.}
  \item $p\neq3$, $p\mid x^2+x+1$: then $\vp_3(x-1)$ is odd, $x-1=3^{2j+1}m^2$. If $j\ge1$ then $27\mid s^3-2$, so $s^3\equiv2\pmod 9$, impossible (cubes mod $9$). So $j=0$, $s^3=3m^2+2$; multiplying by $27$: $(9m)^2=(3s)^3-54$ with $3\mid X=3s$. By Lemma~\ref{lem:mordell} the only integral points have $X=7$, not divisible by $3$. \emph{Closed.}
  \end{itemize}
\end{itemize}

\subsection*{Case 3: $x\equiv2\pmod3$ ($g_-=1$, $g_+=3$)}
Chan side: $x-1=u^2$ ($u$ odd), $x^2+x+1=p^{3+2k}\cdot\square$. She side: $\vp_3(x^2-x+1)=1$ exactly. We first dispose of the case $q=3$:
\begin{itemize}
\item $q=3$: here $x^3+1=9b^3$, so $\vp_3(x^3+1)=2+3\vp_3(b)$, while $\vp_3(x^3+1)=\vp_3(x+1)+1$ by Lemma~\ref{lem:gcd}; hence $\vp_3(x+1)=1+3\vp_3(b)$ and $\vp_3(x^2-x+1)=1$, with all other prime exponents in $x^2-x+1$ being $\equiv0\pmod3$. Thus $x^2-x+1=3s^3$, and Lemma~\ref{lem:she2} gives $x\in\{-1,2\}$; $x=2$ fails since $x^3-1=7$ is not of the form $p^3y^2$. \emph{Closed.}
\item $q\neq3$: from $x^3+1=q^2b^3$ and $\vp_3(x^2-x+1)=1$ we get $\vp_3(x+1)=3\vp_3(b)-1\equiv2\pmod 3$, in particular $9\mid x+1$. Sub-split:
  \begin{itemize}
  \item $q\mid x+1$: then $x^2-x+1=3t^3$ ($\vp_3=1$, $q\nmid$, all other exponents $\equiv0\bmod 3$), so $x\in\{-1,2\}$ by Lemma~\ref{lem:she2}; as before. \emph{Closed.}
  \item $q\mid x^2-x+1$: then $x+1=3^{3\vp_3(b)-1}\cdot s^3=9c^3$ for an integer $c$, and $x-1=u^2$ gives $u^2+2=9c^3$; multiplying by $81$: $(9u)^2=(9c)^3-162$. By Lemma~\ref{lem:mordell}, $Y^2=X^3-162$ has no integral points at all. \emph{Closed.}
  \end{itemize}
\end{itemize}

This completes the case analysis: both families reduce to the branches recorded above, all of which are empty except the two leading to \eqref{eq:wall}.

\section{The equation $t^6+t^3+1=3w^2$}\label{sec:wall}

Write $\Phi_9(t)=t^6+t^3+1$ for the ninth cyclotomic polynomial, so \eqref{eq:wall} reads $\Phi_9(t)=3w^2$. The smooth projective model of
\[
\mathcal{C}:\; Y^2=3t^6+3t^3+3 \qquad (Y=3w)
\]
is a curve of genus $2$. Since the right-hand side is palindromic, $\mathcal{C}$ carries the involution $\sigma:(t,Y)\mapsto(1/t,\,Y/t^3)$ defined over $\Q$, and hence (besides the hyperelliptic involution $\iota$) the second involution $\tau=\sigma\iota$. The quotients $\mathcal{C}/\sigma$ and $\mathcal{C}/\tau$ are elliptic, and $\operatorname{Jac}(\mathcal{C})$ is isogenous over $\Q$ to their product. Concretely:

\begin{lemma}\label{lem:quotients}
Set $u=t+1/t$, and
\[
V=\frac{Y(t+1)}{t^2},\qquad W=\frac{Y(t-1)}{t^2}.
\]
Then on $\mathcal{C}$ one has the exact identities
\[
V^2=3(u+2)(u^3-3u+1)=:q_+(u),\qquad W^2=3(u-2)(u^3-3u+1)=:q_-(u).
\]
\end{lemma}
\begin{proof}
Multiply $Y^2=3(t^6+t^3+1)$ by $(t\pm1)^2/t^4$ and use
\[
(t\pm1)^2\,(t^6+t^3+1)=t^4\,(u\pm2)\,(u^3-3u+1),
\]
which is a polynomial identity in $t$ after substituting $u=t+1/t$ (verified exactly; it follows from $u^3-3u+2=(u-1)^2(u+2)$ and $u^3-3u-2=(u+1)^2(u-2)$ together with $t^3+t^{-3}=u^3-3u$ and $(t\pm1)^2/t=u\pm2$).
\end{proof}

\begin{proof}[Proof of Theorem~\ref{thm:wall}]
Let $(t,w)\in\Q^2$ satisfy \eqref{eq:wall} and put $Y=3w$. Note $t\neq0$ (else $3w^2=1$) and $u=t+1/t\neq-1$ (else $t^2+t+1=0$, no rational root). By Lemma~\ref{lem:quotients}, $(u,W)$ is a rational point on the genus-one quartic curve
\[
\mathcal{D}_-:\; W^2=q_-(u)=3(u-2)(u^3-3u+1).
\]
The quartic $q_-$ is squarefree (its cubic factor has discriminant $81$ and does not vanish at $u=2$), so $\mathcal{D}_-$ is smooth of genus one; its two points at infinity satisfy $W_\infty^2=3$ and are conjugate over $\Q(\sqrt3)$, hence not rational. $\mathcal{D}_-$ has the rational point $(2,0)$.

We claim $\mathcal{D}_-(\Q)=\{(2,0)\}$. Substituting $u=2+1/T$, $W=S/T^2$ transforms $\mathcal{D}_-$ birationally into
\[
S^2=9T^3+27T^2+18T+3,
\]
and the further substitution $(x,y)=(9T+9,\,9S)$ gives the elliptic curve
\[
E:\; y^2=x^3-81x+243,
\]
of conductor $1296$. This is the curve \texttt{1296.k1} in the LMFDB \cite{LMFDB} (Cremona label \texttt{1296f2}): it has Mordell--Weil rank $0$ and trivial torsion, so $E(\Q)=\{\mathcal{O}\}$ --- there is no affine rational point. (We verified rank $0$ and triviality of torsion independently in SageMath with \texttt{rank(proof=True)}.) Running the birational maps backwards, any rational point of $\mathcal{D}_-$ with $u\neq2$ would produce an affine rational point of $E$; hence $u=2$ is the only possibility, and then $W^2=q_-(2)=0$.

Pulling back to $\mathcal{C}$: $u=t+1/t=2$ forces $(t-1)^2=0$, so $t=1$ and $Y^2=9$, $w=\pm1$. \end{proof}

\begin{remark}
The companion quotient $\mathcal{D}_+:V^2=q_+(u)$ has Jacobian $y^2=x^3-189x+999$ of rank $1$ (LMFDB \texttt{1296.b1}, Cremona \texttt{1296l2}); this is why the descent must go through $\mathcal{D}_-$.
\end{remark}

\begin{remark}[no congruence sieve suffices]\label{rem:covering}
Equation \eqref{eq:wall} cannot be resolved by congruences alone, in the following precise sense. Substituting $z=t^3$, any solution gives a point on the conic $z^2+z+1=3w^2$; writing $2z+1=3c$, $2w=b$ reduces the conic to the Pell equation $b^2-3c^2=1$, whose solutions with $b$ even form a single orbit $\{z_k\}_{k\in\Z}$ of the recurrence $z_{k+1}=14z_k-z_{k-1}+6$ through $(z_0,z_1)=(-2,1)$, the sign choice corresponding to the symmetry $z\mapsto-1-z$. Asking whether \eqref{eq:wall} has a solution with $t\neq1$ is asking whether the orbit contains a cube other than $z_1=1$. Since $z_1=1$ \emph{is} a cube, its residue class modulo any period of the sequence is a ``shadow'' that no congruence can eliminate --- and via the symmetry it also protects a class in the negative branch. Hence no finite system of congruence conditions can prove the nonexistence of further cubes in the orbit, and some global (here: elliptic) input is unavoidable. A direct check of the recurrence (in exact integer arithmetic) confirms that no $z_k$ with $k\le2000$ other than $z_1$ is a cube, the largest candidates having more than $2200$ decimal digits.
\end{remark}

\begin{remark}[relation to Nagell--Ljunggren and to powers in recurrences]
Since $\Phi_9(t)=(t^9-1)/(t^3-1)$, equation \eqref{eq:wall} is adjacent to the classical equations $(x^n-1)/(x-1)=3y^m$ of Nagell--Ljunggren type. Note however that Nagell's theorem on $x^2+x+1=3y^n$ ($n>2$) does not apply: in \eqref{eq:wall} the square sits on the right-hand side ($m=2$), while the cube condition sits inside the conic variable $z=t^3$. In the language of Remark~\ref{rem:covering}, \eqref{eq:wall} asks for the cubes in a fixed binary recurrence orbit; the effective methods underlying the finiteness results of Peth\H{o} \cite{Petho} and Shorey--Stewart \cite{ShoreyStewart} guarantee that there are only finitely many such cubes, but do not by themselves produce the explicit solution set. For the closely related factorisation $(2z+1)(z^2+z+1)=y^\ell$ see Bennett, Patel, and Siksek \cite{BPS}. We are not aware of a published complete determination of the integer (or rational) solutions of \eqref{eq:wall}; the contribution here is the explicit resolution via the rank-$0$ elliptic quotient. 
\end{remark}

\section{A modular refinement of Beckon's condition}\label{sec:beckon}

Beckon \cite{Beckon} observed that the smallest member $n$ of a triple of consecutive powerful numbers must satisfy $n\equiv7,27,35\pmod{36}$; this elementary but useful constraint appears not to have been taken up in the recent literature. 
The underlying mechanism generalizes to any prime-square modulus.

\begin{proof}[Proof of Proposition~\ref{prop:beckon}]
If $m$ is powerful and $\ell$ prime, then $\vp_\ell(m)\neq1$; in particular, if $m\equiv c\pmod{\ell^2}$ with $\ell\mid c$, $\ell^2\nmid c$, then $m$ is not powerful. Call $r\bmod \ell^2$ \emph{admissible} if none of $r,r+1,r+2$ is $\equiv\ell c'\pmod {\ell^2}$ with $\ell\nmid c'$. For $\ell\ge5$ the $\ell-1$ forbidden residues $\ell c'$ are pairwise at distance $\ge5$, so the $3(\ell-1)$ shifted prohibitions are distinct, leaving $\ell^2-3(\ell-1)=\ell^2-3\ell+3$ admissible classes. For $\ell=2$: $n$ must be odd and $n+1\equiv0\pmod4$, leaving the class $3\bmod4$ ($1$ of $4$). For $\ell=3$: the windows avoiding $\{3,6\}\bmod 9$ start at $r\in\{0,7,8\}$ ($3$ of $9$). By CRT the admissible classes multiply: modulo $36$ one obtains exactly Beckon's $\{7,27,35\}$; modulo $900=2^2 3^2 5^2$ there are $3\cdot13=39$ classes; modulo $44100=2^23^25^27^2$ there are $39\cdot31=1209$. (The explicit lists were computed by two independent implementations with matching output.) Finally, $\prod_{5\le\ell\le L}\bigl(1-\tfrac{3}{\ell}+\tfrac{3}{\ell^2}\bigr)\asymp(\log L)^{-3}$ by Mertens' theorem, which proves (iii).
\end{proof}

\section{Concluding remarks}

Together with \cite{Chan2025} and \cite{She2025}, Theorem~\ref{thm:main} completes the four shape-combinations $\{p^3\cdot\square,\,p^2\cdot\mathrm{cube}\}^2$ for consecutive powerful triples around a cube in which both outer members are divisible by a single controlling prime to an exact pattern. Natural next steps in this program include the relaxations $p=q$ and composite cofactors, and She's question whether $x^n-1$ can be powerful for $x>1$, $n>2$. The bielliptic reduction of Section~\ref{sec:wall} may be of independent use for other palindromic sextic obstructions in this circle of problems.

\paragraph{Acknowledgements.}
Substantial parts of this work, including the reduction of Section~\ref{sec:wall}, were developed with the assistance of Claude (Fable~5 model, Anthropic). All steps were independently verified as described above: SageMath (\texttt{integral\_points}, \texttt{rank} with \texttt{proof=True}, in two independent runs), exact symbolic identities in SymPy, and cross-checks against the LMFDB. The author reviewed all arguments and takes full responsibility for the content.

\begin{thebibliography}{99}

\bibitem{Beckon} E.~Beckon, \emph{On consecutive triples of powerful numbers}, Rose--Hulman Undergrad.\ Math.\ J. \textbf{20} (2019), no.~2, Article~3.

\bibitem{BennettGhadermarzi} M.~A.~Bennett and A.~Ghadermarzi, \emph{Mordell's equation: a classical approach}, LMS J.\ Comput.\ Math. \textbf{18} (2015), 633--646.

\bibitem{BPS} M.~A.~Bennett, V.~Patel, and S.~Siksek, \emph{Perfect powers that are sums of consecutive cubes}, Mathematika \textbf{63} (2017), no.~1, 230--249; arXiv:1603.08901.

\bibitem{Chan2025} T.~H.~Chan, \emph{A note on three consecutive powerful numbers}, Integers \textbf{25} (2025), \#A7; arXiv:2503.21485.

\bibitem{Erdos76} P.~Erd\H{o}s, \emph{Problems and results on consecutive integers}, Publ.\ Math.\ Debrecen \textbf{23} (1976), no.~3--4, 271--282.

\bibitem{GPZ} J.~Gebel, A.~Peth\H{o}, and H.~G.~Zimmer, \emph{On Mordell's equation}, Compositio Math. \textbf{110} (1998), 335--367.

\bibitem{Golomb} S.~W.~Golomb, \emph{Powerful numbers}, Amer.\ Math.\ Monthly \textbf{77} (1970), 848--852.

\bibitem{LMFDB} The LMFDB Collaboration, \emph{The L-functions and modular forms database}, \url{https://www.lmfdb.org}, 2026. Elliptic curves \texttt{1296.k1} and \texttt{1296.b1}.

\bibitem{MollinWalsh86a} R.~A.~Mollin and P.~G.~Walsh, \emph{On powerful numbers}, Internat.\ J.\ Math.\ Math.\ Sci. \textbf{9} (1986), 801--806.

\bibitem{MollinWalsh86b} R.~A.~Mollin and P.~G.~Walsh, \emph{A note on powerful numbers, quadratic fields and the Pellian}, C.~R.\ Math.\ Rep.\ Acad.\ Sci.\ Canada \textbf{8} (1986), 109--114.

\bibitem{Petho} A.~Peth\H{o}, \emph{Perfect powers in second order linear recurrences}, J.\ Number Theory \textbf{15} (1982), no.~1, 5--13.

\bibitem{She2025} J.~She, \emph{Nonexistence of consecutive powerful triplets around cubes with prime-square factors}, Integers \textbf{25} (2025), \#A103; arXiv:2507.16828v3.

\bibitem{ShoreyStewart} T.~N.~Shorey and C.~L.~Stewart, \emph{On the Diophantine equation $ax^{2t}+bx^ty+cy^2=d$ and pure powers in recurrence sequences}, Math.\ Scand. \textbf{52} (1983), no.~1, 24--36.

\bibitem{Tzanakis} N.~Tzanakis, \emph{The diophantine equation $x^3-3xy^2-y^3=1$ and related equations}, J.\ Number Theory \textbf{18} (1984), no.~2, 192--205.

\bibitem{vanDoorn} W.~van Doorn, \emph{Three-term arithmetic progressions of consecutive powerful numbers}, arXiv:2605.06697 (2026).

\end{thebibliography}

\end{document}
